%==============================================================================
\begin{abstract}
On-chip slack sensors used in Silicon Lifecycle Management (SLM) report a
quantity that mixes a permanent, slowly-drifting aging component with a
reversible component driven by process, voltage and temperature (PVT). Prior
work quantifies this contamination with the coefficient of determination
($R^2$) of a linear regression of slack on temperature and voltage---a measure
that captures only linear dependence and yields no transferable unit. This work
places an \emph{information-theoretic measurement layer} on top of an existing
low-cost accelerated-aging burn-in platform and its synchronized
\mbox{(slack, $\Tdie$, $\Vcore$, time)} logs. Treating the sensor as a noisy
communication channel, we (i)~replace $R^2$ with the mutual information
$\MI{\slk}{\Tdie,\Vcore}$ as the correct, non-linear measure of environmental
contamination; (ii)~quantify, in bits, how much uncertainty PVT logging removes
via $\Ent{\slk}-\Ent{\slk\mid\Tdie,\Vcore}$; (iii)~estimate the sensor's
effective resolving power through a channel-capacity heuristic and a Binary
Symmetric Channel model of the aging-alarm decision; (iv)~derive a
noise-limited detectability bound for remaining-useful-life (RUL) estimation,
via both a Cram\'er--Rao/Bayesian point estimate and a sequential Wiener-process/
Kalman-filter treatment that yields a denoised aging trajectory with credible
bands and a full, closed-form Inverse-Gaussian RUL distribution. We further
apply MDL model selection and ICA-based blind source separation to the same
data, and a KL/Jensen--Shannon comparison of two validation benches. On a
\valDurationHours~h PID-controlled campaign (\valPostSoakSamples{} post-soak
samples; effective sample size $\neff=\valNeff$) we measure
$\Ent{\slk}=\valHs$~bits and $\MI{\slk}{\Tdie,\Vcore}=\valMIstv$~bits
(KSG cross-check \valKSGstv~bits). This measured dependence exceeds the
Gaussian-equivalent MI of the prior work's linear fit
($R^2=\valRsqLinear\Rightarrow\valMIlinEquiv$~bits) by \valMIexcessPct\,\%,
directly evidencing the non-linear environmental coupling that $R^2$ cannot
capture. The contribution is a correct, transferable noise metric for SLM
sensors and a principled, optimizable definition of burn-in quality as the
minimization of environment--sensor mutual information.
\end{abstract}

\noindent\textbf{Keywords:} information theory; mutual information; entropy;
channel capacity; silicon lifecycle management; aging sensor; slack;
burn-in; remaining useful life.
